\section{Experimental Setup}
\label{sec:experiments}

The preceding hourly and monthly experiments are retained as evidence from the
two source manuscripts. They are not treated as a matched comparison. The
extension adds only experiments that align resolution, context, horizon,
origins, and evaluation rules for TimesNet and DilatedRNN.

\begin{table}[htbp]
\centering
\caption{Existing evidence and experiments required by the journal extension.}
\label{tab:study_design}
\small
\begin{tabular}{p{0.13\textwidth}p{0.15\textwidth}p{0.17\textwidth}p{0.16\textwidth}p{0.23\textwidth}}
\toprule
\textbf{Resolution} & \textbf{Model} & \textbf{Context} & \textbf{Horizon} & \textbf{Status} \\
\midrule
Hourly & TimesNet & 336 h & 24/48/72 h & Complete source experiment \\
Hourly & DilatedRNN & To be matched & 24/48/72 h & Required \\
Daily & TimesNet & To be fixed & Multiple days & Required \\
Daily & DilatedRNN & To be matched & Multiple days & Required \\
Monthly & DilatedRNN & 12 months & 1 month & Complete source experiment \\
Monthly & TimesNet & 12 months & 1 month & Required \\
Multi-month & Both & To be matched & To be fixed & Required only if defined as long-term \\
\bottomrule
\end{tabular}
\end{table}



\subsection{Hourly protocol}

The available TimesNet protocol uses 336 hourly observations as context and
predicts 72 hours directly. The 24- and 48-hour results are prefixes of the
same 72-hour output. Training spans 2019--2022, 2023 is used for model
selection, the final model is refitted on 2019--2023, and 2024 is reserved for
retrospective testing.

\subsection{Monthly protocol}

The available DilatedRNN protocol uses a 12-month context and predicts the next
monthly mean GHI. It evaluates the 12 chronological origins of 2024 after
training and model selection on pre-test observations. This is a one-step
monthly task; it will not be described as a multi-month forecast.

\subsection{Daily and matched long-horizon protocols}

\DraftNote{Execute the daily TimesNet and DilatedRNN experiments at identical
origins. Fix the daily horizons before inspecting test performance. If
``long-term'' means multiple monthly steps, execute both architectures with the
same direct or recursive forecast strategy and report each horizon separately.}

The primary configuration table will contain only the parameters needed to
understand and reproduce the comparison. Complete search spaces, selected
epochs by seed, and secondary settings are assigned to
Appendix~\ref{app:hyperparameters}.

\begin{table}[htbp]
\centering
\caption{Verified configurations in the two source experiments.}
\label{tab:hyperparameters}
\small
\begin{tabular}{p{0.20\textwidth}p{0.31\textwidth}p{0.34\textwidth}}
\toprule
\textbf{Item} & \textbf{Hourly TimesNet} & \textbf{Monthly DilatedRNN} \\
\midrule
Context/output & 336 h / 72 h & 12 months / 1 month \\
Main width & $d_{\mathrm{model}}=8$, $d_{\mathrm{ff}}=16$ & Three recurrent layers, 16 units each \\
Temporal blocks & One TimesBlock, $k=3$ & Dilations 1--2--4 \\
Optimizer & Adam, learning rate $10^{-3}$ & Adam, learning rate $10^{-3}$ \\
Batch size & 128 & 32 \\
Loss & MSE & MSE \\
Seeds & 42 & 11, 23, 42, 67, 89 \\
\bottomrule
\end{tabular}
\end{table}



